\documentstyle[% 


righttag% 


,11pt% 


,amssymb% 


,russian%               remove in Eng. 


,izvru%                 Set izven% in Eng. 


%,izvru-ks             for brief communications 


]{amsart} 


 


%  


\newcommand{\Grad}{\operatorname{grad}} 


\newcommand{\Div}{\operatorname{div}} 


\newcommand{\const}{\operatorname{const}} 


\newcommand{\im}{\operatorname{Im}} 


\newcommand{\tts}[1]{} 


 


 


%\hoffset=-15mm%                  For Eng. 


\setcounter{page}{43} 


\god{2004} 


\nomer{2~(501)} %                        Remove in Eng. 


%\nomer{Vol.\,48, No.\,2}%               Set in Eng. 


%\str{\pageref{begin}--\pageref{end}}%  Set in Eng. 


\udk{517.946} 


 


\author{\it О.И.\,Махмудов} 


% NB:   ^^ remove in Eng. 


\title{Задача Коши для системы уравнений теории упругости\\ 


и термоупругости в пространстве} 


\thanks{} 


 


 


\date{} 


%\pagestyle{myheadings}%         Set in Eng. 


\pagestyle{plain} %              Remove in Eng. 


\begin{document} 


%\thispagestyle{plain}%          Set in Eng. 


\maketitle 


\label{begin} 


 


 


Система уравнений теории упругости является предметом многих 


исследований. В математической физике для этой системы ставятся три 


основные задачи, которые отличаются друг от друга видом краевых условий. 


В первой задаче на границе области задаются смещения, во второй задаче --- 


напряжения, в третьей (смешанной задаче) на части границы задаются 


напряжения, а на оставшейся части --- смещения. 


 


В данной работе рассматривается задача аналитического 


продолжения решения системы уравнений теории упругости в 


пространственной неограниченной области по его значениям и значениям его 


напряжения на части границы этой области, т.\,е. задача Коши. 


 


Система уравнений теории упругости эллиптическая, задача Коши для 


эллиптических уравнений неустойчива относительно малого изменения 


данных, т.\,е. некорректна (пример Адамара [1]). В некорректных задачах 


теорема существования не доказывается, существование предполагается 


заданным априори. Более того, предполагается, что решение принадлежит 


некоторому заданному подмножеству функционального пространства, 


обычно компактному, ([2], с.\,4). Единственность решения следует из общей 


теоремы Холмгрена ([3], с.\,58). 


 


После установления единственности в теоретических исследованиях 


некорректных задач возникают важные вопросы получения оценки условной 


устойчивости и построения регуляризирующих операторов. В 1926~г. 


Т.\,Карлеман ([2], с.\,41) построил 


формулу, которая связывает значения аналитической 


функции комплексного переменного в точках области с ее значениями на 


куске границы этой области. На основе этой формулы в ([2], с.\,34) 


введено понятие функции Карлемана задачи Коши для уравнения Лапласа и в 


некоторых случаях указан способ ее построения. Конструкция функции 


Карлемана дает возможность в этих задачах построить регуляризацию и 


получить оценку условной устойчивости. 


 


На протяжении последних десятилетий не ослабевал интерес к 


классической некорректной задаче математической физики. Это направление 


в исследовании свойств решений задачи Коши для уравнения Лапласа начато 


в 50-х годах в работах [2], [4]--[6] 


и развивалось впоследствии в [7]--[14]. 


\medskip 


 


{\bf 1.} Пусть $x=(x_1,x_2,x_3)$, $y=(y_1,y_2,y_3)$ --- точки трехмерного 


евклидова 


пространства $E^3$ и упругая среда $D$ есть неограниченная односвязная 


область 


в $E^3$ с кусочно-гладкой границей, состоящей из плоскости $\Sigma:y_3=0$ и 


гладкой поверхности $S$, лежащей в полупространстве $y_3>0$, 


т.\,е. $\partial D=SU\Sigma$. Рассмотрим в области $D$ системы уравнений 


теории 


упругости в векторной форме ([15], с.\,10) 


\begin{equation} 


\mu\Delta u(y)+(\lambda+\mu)\Grad\Div u(y)=0; 


\end{equation} 


здесь $u=(u_1;u_2;u_3)$ --- вектор смещения, $\Delta$ --- оператор Лапласа, 


$\lambda$, $\mu$ --- постоянные 


Ламе рассматриваемой упругой среды. 


\medskip 


 


{\it Постановка задачи}. Известны данные Коши решения системы (1) на 


поверхности $S$: 


\begin{equation} 


\begin{split} 


u(y)&=f(y),\quad y\in S,\\ 


T(\partial_y,n)u(y)&=g(y),\quad y\in S, 


\end{split} 


\end{equation} 


где 


$$ 


T(\partial_y,n)=\|T_{kj}(\partial_y,n)\|_{3\times 3}=\bigg\| 


\lambda n_k\frac{\partial}{\partial y_j}+\mu n_j 


\frac{\partial}{\partial y_k}+\mu\delta_{kj} 


\frac{\partial}{\partial n}\bigg\|_{3\times 3}


$$ 


--- оператор напряжения, 


$n=(n_1,n_2,n_3)$ --- единичный вектор нормали к поверхности $S$, 


$f=(f_1,f_2,f_3)$, $g=(g_1,g_2,g_3)$ --- заданные непрерывные вектор-функции 


на 


$S$, $\delta_{kj}$ --- символ Кронекера. 


 


Требуется определить функцию $u(y)$ в $D$, исходя из заданных $f$ и $g$, 


т.\,е. решить задачу аналитического продолжения решения системы уравнений 


теории упругости в пространственной неограниченной области по ее 


значениям $f$ и значениям ее напряжений $g$ на гладком куске $S$ границы. 


 


Пусть вместо $f(y)$ и $g(y)$ заданы их приближения $f_\delta$ и $g_\delta$ с 


точностью 


$\delta\in(0,1)$ (в метрике $C$). В данной работе строится семейство 


вектор-функций 


$U(x,f_\delta,g_\delta)=U_{\sigma\delta}$, зависящее от параметра $\sigma$, и 


доказывается, что при 


некоторых условиях и специальном выборе параметра $\sigma(\delta)$ при 


$\delta\to0$ 


семейство $U_{\sigma\delta}(x)$ сходится в обычном смысле к решению задачи 


(1)--(2). 


 


Следуя [16], функцию $U_{\sigma\delta}(x)$ назовем 


регуляризованным решением задачи Коши для системы теории упругости. 


Регуляризованное решение определяет устойчивость метода приближенного 


решения задачи. 


 


С использованием результатов работ [2], [7], [8] по задаче Коши для 


уравнения Лапласа в данной статье построена матрица Карлемана в явном виде 


и на ее основе --- регуляризованное решение задачи Коши для систем уравнений 


теории упругости. В [17] приведены теоремы существования 


матрицы Карлемана и критерий разрешимости более широкого класса 


краевых задач для эллиптических систем. 


Ранее в [9], [17] доказано, что матрица Карлемана существует во 


всякой задаче Коши для решений эллиптических систем, если только данные 


Коши задаются на граничном множестве положительной меры. Поскольку 


в данной статье идет речь о явных формулах, то построение матрицы 


Карлемана в элементарных и специальных функциях представляет 


значительный интерес. 


 


Функция Карлемана задачи Коши для уравнения Лапласа и ей 


близких в случае, когда $\partial D\setminus S$ --- часть поверхности 


конуса, построена 


в [7], [8]. Матрицу Карлемана задачи Коши для уравнения 


Коши--Римана в случае, когда $S$ --- произвольное множество положительной 


меры, построил Л.А.\,Айзенберг [9]. В развитие идеи С.Е.\,Мергеляна [5], 


указавшего способ построения функции Карлемана задачи Коши для 


уравнения Лапласа в случае, когда $S$ --- кусок с гладким краем границы 


односвязной области, на основе теорем об аппроксимации 


в [17] построена матрица Карлемана для эллиптических систем 


(см. также [10]--[14]). 


В [10]--[14] для плоских ограниченных и неограниченных 


областей и для пространственных ограниченных областей построена матрица 


Карлемaна в явном виде и на ее основе -- регуляризованное решение задачи 


Коши. 


 


Введем матричный дифференциальный оператор 


$A(\partial_x)=\|A_{ij}(\partial_x)\|_{3\times 3}$, где 


$A_{ij}(\partial_x)=\delta_{ij}\mu\Delta+(\lambda+\mu) 


\frac{\partial^2}{\partial x_i\partial x_j}$. 


Тогда уравнение равновесия однородной изотропной упругой среды (1) 


запишется в матричной форме $A=(\partial_x)u=0$. 


 


 


\begin{defn} 


Матрица $\Gamma(y,x)=\|\Gamma_{kj}(y,x)\|_{3\times 3}$ называется матрицей 


фундаментальных решений ([15], с.\,32), где 


\begin{gather*} 


\Gamma_{kj}(y,x)=\frac{\lambda'\delta_{kj}}{|y-x|}+\mu' 


\frac{(y_k-x_k)(y_j-x_j)}{|y-x|^3},\quad k,j=1,2,3,\\ 


% 


|y-x|=\sqrt{(y_1-x_1)^2+(y_2-x_2)^2+(y_3-x_3)^2},\\ 


% 


\lambda'=(\lambda+3\mu)[4\pi\mu(\lambda+2\mu)]^{-1},\quad 


\mu'=(\lambda+\mu)[4\pi\mu(\lambda+2\mu)]^{-1}. 


\end{gather*} 


 


Матрица $\Gamma(y,x)$ симметричная, и каждые ее столбец и строка 


удовлетворяют уравнению упругости в произвольной точке $x\in E^3$, кроме 


$x=y$. Таким образом, $A(\partial_x)\Gamma(y,x)=0$, $x\ne y$. 


Здесь нуль в правой части обозначает матрицу размера $3\times3$ с нулевыми 


элементами. 


\end{defn} 


 


 


Следуя [2], приведем 


 


 


\begin{defn} 


Матрицей Карлемана задачи (1)--(2) называется $3\times3$-матрица 


$\Pi(y,x,\sigma)$, 


удовлетворяющая следующим двум условиям: 


 


1) $\Pi(y,x,\sigma)=\Gamma(y,x)+G(y,x,\sigma)$, 


где $\sigma$ --- положительный числовой параметр, матрица $G(y,x,\sigma)$ по 


переменной $y$ удовлетворяет системе (1) всюду в области $D$, $\Gamma(y,x)$ 


--- 


матрица фундаментальных решений уравнений (1); 


 


2) $\int\limits_\Sigma(|\Pi(y,x,\sigma)|+ |T(\partial_y,n) \Pi(y,x,\sigma)|) 


ds_y\le \varepsilon(\sigma)$ при фиксированном $x\in D$, 


где $\varepsilon(\sigma)\to0$ при $\sigma\to\infty$; здесь и далее $|\Pi|$ 


означает евклидову норму 


матрицы $\Pi=\|\Pi_{kj}\|$, т.\,е. 


$|\Pi|=\Big(\sum\limits^3_{k,j=1}\Pi^2_{kj}\Big)^{1/2}$, в частности 


$|u|=\Big(\sum\limits^3_{k=1}u^2_k\Big)^{1/2}$ для 


вектора $u$. 


\end{defn} 


 


 


\begin{defn} 


Вектор-функция $u(y)=(u_1(y),u_2(y),u_3(y))$ 


называется регулярной в $D$, если $u_j(y)\in C^2(D)\cap C^1(\ov D)$, 


$j=1,2,3$, т.\,е. если 


она непрерывна вместе со своими частными производными второго порядка 


в $D$ и первого порядка на $\ov D=D\cup\partial D$ (если $D$ --- 


неограниченная область, то 


непрерывность $u(y)$ и ее частных производных требуется лишь в конечных 


точках $\partial D$). 


\end{defn} 


 


 


В теории уравнений в частных производных важную роль играют 


представления решений этих уравнений в виде функции типа потенциала. 


Приведем здесь формулу Сомилиана ([15], с.\,57). 


 


 


\begin{thm} 


Всякое регулярное решение $u(y)$ уравнения $(1)$ в области 


$D$ определяется формулой 


$$ 


u(x)=\int_{\partial D}[\Gamma(y,x)\{T(\partial_y,n)u(y)\}- 


\{T(\partial_y,n)\Gamma(y,x)\}'u(y)]ds_y,\quad x\in D. 


$$ 


\end{thm} 


 


 


Поскольку матрица Карлемана отличается от матрицы 


фундаментальных решений на решение транспонированной системы, то 


формула Сомилиана остается справедливой, если в ней заменить 


фундаментальное решение на матрицу Карлемана. 


 


Пусть область $D\subset E^3$ лежит внутри слоя $0<y_3<h$, 


$h=\frac{\pi}{\rho}$, $\rho>0$, 


граница ее состоит из плоскости $\Sigma:y_3=0$ и гладкой поверхности $S$, 


заданной уравнением $y_3=F(y_1,y_2)$, удовлетворяющей условиям 


$0<F(y_1,y_2)\le h$, $|\Grad F(y_1,y_2)|\le\const<\infty$, $y_1,y_2\in E^2$. 


Обозначим через $A(D)$ пространство всех регулярных в $D$ решений 


системы (1), а через $A_\rho(D)$ --- класс функций из $A(D)$, 


удовлетворяющих следующему условию роста: 


$$ 


A_\rho(D)=\{u(y)\in A(D):|u(y)|+|\Grad u(y)|\le 


\exp(o\exp\rho|y_1|),\ \; y\to\infty,\ \; \rho>0\}. 


$$ 


 


С целью построения приближенного решения задачи (1)--(2) 


рассмотрим матрицу 


\begin{equation} 


\Pi(y,x,\sigma)=\bigg\|\lambda'\delta_{kj}\Phi(y,x,\sigma)- 


\mu'(y_j-x_j)\frac{\partial\Phi(y,x,\sigma)}{\partial y_k} 


\bigg\|_{3\times 3}, 


\end{equation} 


где 


\begin{gather} 


\lambda'=\frac{\lambda\mu+3\mu}{4\pi\mu(\lambda+2\mu)},\quad 


\mu'=\frac{\lambda+\mu}{4\pi\mu(\lambda+2\mu)},\notag \\ 


% 


\Phi(y,x,\sigma)=[-\pi K(x_3)]^{-1}\int^\infty_0 \im 


\frac{K(\omega)}{\omega-x_3}\cdot\frac{du}{\sqrt{u^2+\alpha^2}}, \\ 


% 


K(\omega)=(\omega-x_3+2h)^{-1}\exp (\sigma\cdot\omega),\quad 


\omega=i\sqrt{u^2+\alpha^2}+y_3,\notag \\ 


% 


\alpha^2=(y_1-x_1)^2+(y_2-x_2)^2,\quad 0<x_3<h,\quad 


K(x_3)=(2h)^{-1}\exp(\sigma x_3).\notag 


\end{gather} 


 


В [7] доказана 


 


 


\begin{lem} 


Функция $\Phi(y,x,\sigma)$, определенная формулой $(4)$, представима 


в виде 


$$ 


\Phi(y,x,\sigma)=\frac{1}{r}+\psi(y,x,\sigma),\quad 


r=|y-x|, 


$$ 


где $\psi(y,x,\sigma)$ --- некоторая функция, определенная для всех значений 


$y$, $x$ и 


гармоническая по переменной $y$ во всем $E^3$. 


\end{lem} 


 


 


\begin{lem} 


Матрица $\Pi(y,x,\sigma)$, определенная формулой $(3)$, является 


матрицей Карлемана задачи $(1)$--$(2)$, т.\,е. представима в виде 


\begin{equation} 


\Pi(y,x,\sigma)=\Gamma(y,x)+G(y,x,\sigma), 


\end{equation} 


где $G(y,x,\sigma)=\|G_{kj}(y,x,\sigma)\|_{3\times 3}$ --- матрица, 


определенная для всех значений 


$y$, $x$, и по переменной $y$ удовлетворяющая системе $(1)$ во всем $E^3$,


\begin{equation} 


A(\partial_y)G(y,x,\sigma)=0. 


\end{equation} 


\end{lem} 


 


 


\begin{pf} 


В силу леммы 1 и определения 1 имеем 


\begin{gather*} 


\Pi(y,x,\sigma)=\bigg\|\lambda'\delta_{kj}\bigg[\frac{1}{r}+ 


\psi(y,x,\sigma)\bigg]-\mu'(y_j-x_j)\frac{\partial}{\partial y_k} 


\bigg[\frac{1}{r}+\psi(y,x,\sigma)\bigg]\bigg\|_{3\times 3}= \\ 


% 


=\bigg\|\frac{\lambda'\delta_{kj}}{|y-x|}+\mu' 


\frac{(y_k-x_k)(y_j-x_j)}{|y-x|^3}\bigg\|_{3\times 3}+ \\ +\bigg\| 


\lambda'\delta_{kj}\psi(y,x,\sigma)-\mu'(y_j-x_j) 


\frac{\partial\psi(y,x,\sigma)}{\partial y_k}\bigg\|_{3\times 3}=


\Gamma(y,x)+G(y,x,\sigma), 


\end{gather*} 


где $\psi(y,x,\sigma)$ --- функция, гармоническая по $y$ во всем $E^3$. 


Отсюда следует (5). Докажем, что матрица 


$$ 


G(y,x,\sigma)=\|G_{kj}(y,x,\sigma)\|_{3\times 3}=\bigg\|\lambda' 


\delta_{kj}\psi(y,x,\sigma)-\mu'(y_j-x_j) 


\frac{\partial\psi(y,x,\sigma)}{\partial y_k}\bigg\|_{3\times 3} 


$$ 


есть регулярное решение уравнения (1) по переменной $y$, включая и точку 


$y=x$. Поскольку 


\begin{gather*} 


\Delta\psi(y,x,\sigma)=0,\quad \Delta=\frac{\partial^2}{\partial y^2_1}+ 


\frac{\partial}{\partial y^2_2}+\frac{\partial^2}{\partial y^2_3},\\ 


% 


\Div G^j(y,x,\sigma)=\frac{1}{2\pi(\lambda+\mu)} 


\frac{\partial}{\partial y_j}\psi(y,x,\sigma), 


\end{gather*} 


то для $k$-й компоненты вектора $A(\partial_y)G^j(y,x,\sigma)$ получим 


\begin{gather*} 


\sum^3_{i=1}A_{kj}(\partial_y)G_{ij}(y,x,\sigma)=\mu\Delta\bigg[ 


\lambda'\delta_{kj}\psi(y,x,\sigma)-\mu'(y_j-x_j) 


\frac{\partial\psi(y,x,\sigma)}{\partial y_k}\bigg]+ \\ 


% 


+(\lambda+\mu)\frac{\partial}{\partial y_k}\Div G^j(y,x,\sigma)= 


-\frac{\lambda+\mu}{2\pi(\lambda+2\mu)} 


\frac{\partial^2\psi(y,x,\sigma)}{\partial y^2_j}+ 


\frac{\lambda+\mu}{2\pi(\lambda+2\mu)} 


\frac{\partial^2\psi(y,x,\sigma)}{\partial y^2_j}=0. 


\end{gather*} 


Таким образом, каждый столбец матрицы $G(y,x,\sigma)$ удовлетворяет 


уравнению (1), т.\,е. имеет место (6) во всем $E^3$. 


 


Из формулы (4) видно, что на $\Sigma(y_3=0)$ функция $\Phi(y,x,\sigma)$, ее 


градиент 


$\nabla\Phi(y,x,\sigma)$ и вторые частные производные 


$\frac{\partial^2\Phi(y,x,\sigma)}{\partial y_i\partial y_j}$ ($i,j=1,2,3$) 


при $\sigma\to\infty$ экспоненциально стремятся к нулю при всех $y_1$, $y_2$


и $x\in E^3$, $x_3>0$. Тогда в силу (3), (4) матрица $G(y,x,\sigma)$ и ее 


напряжение $T(\partial_y,n)\Pi(y,x,\sigma)$ при $\sigma\to\infty$ также 


стремятся к нулю при всех 


$y_1$, $y_2$ и $x\in E^3$, $x_3>0$. В силу определения 2 матрица 


$\Pi(y,x,\sigma)$, 


определенная формулами (3), (4), является матрицей Карлемана для области 


$D$ и части $\Sigma$. 


\end{pf} 


 


 


Предположим, что $u(y)\in A(D)$ ограничена вместе с нормальной 


производной на $\partial D$: 


\begin{equation} 


|u(y)|+|T(\partial_y,n)u(y)|\le M,\quad y\in\Sigma. 


\end{equation} 


В этих предположениях верна формула Сомилиана ([15], с.\,57) 


\begin{equation} 


u(x)=\int_{\partial D}[\Pi(y,x,\sigma)\{T(\partial_y,n)u(y)\}- 


u(y)\{T(\partial_y,n)\Pi(y,x,\sigma)\}]ds_y,\quad x\in D. 


\end{equation} 


Обозначим 


\begin{equation} 


u_\sigma(x)=\int_S[\Pi(y,x,\sigma)\{T(\partial_y,n)u(y)\}- 


u(y)\{T(\partial_y,n)\Pi(y,x,\sigma)\}]ds_y,\quad x\in D. 


\end{equation} 


 


 


\begin{thm} 


Пусть $u(x)\in A_\rho(D)$ удовлетворяет граничному условию $(7)$. 


Тогда 


\begin{equation} 


|u(x)-u_\sigma(x)|\le M C_\rho(x)\sigma^2\exp(-\sigma x_3), 


\quad x\in D, 


\end{equation} 


где $C_\rho(x)=C(\rho)\int\limits_\Sigma|(\sigma r)^{-1}+r^{-2}|ds_y$, 


$r=|y-x|$. 


\end{thm} 


 


 


\begin{pf} 


Согласно формуле (8) имеем 


$$ 


u(x)-u_\sigma(x)=\int_\Sigma[\Pi(y,x,\sigma)\{T(\partial_y,n)u(y)\}- 


u(y)\{T(\partial_y,n)\Pi(y,x,\sigma)\}]ds_y, 


$$ 


а в силу (7) 


\begin{equation} 


|u(x)-u_\sigma(x)|\le M\int_\Sigma|\Pi(y,x,\sigma)|+ 


|T(\partial_y,n)\Pi(y,x,\sigma)|ds_y. 


\end{equation} 


Оценим $\Pi(y,x,\sigma)$ и $|T(\partial_y,n)\Pi(y,x,\sigma)|$. Для этого 


необходимо оценить 


$$ 


\Phi,\ \ \frac{\partial\Phi}{\partial y_i}, 


\ \ \frac{\partial^2\Phi}{\partial y_i\partial y_j},\quad i,j=1,2,3. 


$$ 


По построению 


$$ 


\Phi(y,x,\sigma)=-\frac{2h}{\pi}\exp\sigma(y_3-x_3)\bigg( 


\int^\infty_0\frac{(y_3-x_3)}{u^2+r^2} 


\frac{\sin\sigma\sqrt{u^2+\alpha^2}}{\sqrt{u^2+\alpha^2}}du+ 


\int^\infty_0\frac{1}{u^2+r^2}\cos\sigma\sqrt{u^2+\alpha^2}du\bigg). 


$$ 


Тогда 


\begin{equation} 


|\Phi(y,x,\sigma)|\le C_1(\rho)r^{-1}\sigma\exp\sigma(y_3-x_3). 


\end{equation} 


Нетрудно проверить, что при $i,j=1,2,3$ 


\begin{align} 


\bigg|\frac{\partial\Phi}{\partial y_i}\bigg|&\le 


C_2(\rho)r^{-2}\sigma\exp\sigma(y_3-x_3),\\ 


% 


\bigg|\frac{\partial^2\Phi}{\partial y_i\partial y_j}\bigg|&\le 


C_3(\rho)r^{-3}\sigma^2\exp\sigma(y_3-x_3). 


\end{align} 


Тогда по определению $\Pi(y,x,\sigma)$ и из неравенств (12)--(14) получим 


\begin{align} 


|\Pi(y,x,\sigma)|&\le C_4(\rho)r^{-1}\sigma\exp\sigma(y_3-x_3),\\ 


% 


|T(\partial_y,n)\Pi(y,x,\sigma)|&\le C_5(\rho)r^{-2}\sigma^2 


\exp\sigma(y_3-x_3). 


\end{align} 


Следовательно, из последних неравенств и (11) имеем 


$$ 


|u(x)-u_\sigma(x)|\le M\sigma^2\exp(-\sigma x_3) 


\int_\Sigma[C_4(\rho)r^{-1}+C_5(\rho)r^{-2}]ds_y. 


$$ 


Обозначив 


$C_\rho(x)=\int\limits_\Sigma[C_4(\rho)r^{-1}+C_5(\rho)r^{-2}]ds_y$, получим 


(10). 


\end{pf} 


 


 


Приведем оценку устойчивости. 


 


 


\begin{thm} 


Пусть $u(x)\in A_\rho(D)$ на гиперплоскости $y_3=0$ 


удовлетворяет граничному условию $(7)$, а на $S$ --- условию 


$$ 


|u(y)|+|T(\partial_y,n)u(y)|\le\delta,\quad y\in S,\quad 0<\delta<1. 


$$ 


Тогда 


$$ 


|u(x)|\le C_\rho(x)\frac{2M^{1-x_3/h}}{h^2}\delta^{x_3/h}\ln\bigg( 


\frac{M}{\delta}\bigg)^2,\quad x\in D, 


$$ 


где 


$$ 


C_\rho(x)=C(\rho)\int_{\partial D}[r^{-1}+r^{-2}]ds_y. 


$$ 


\end{thm} 


 


 


{\bf Доказательство.} Из теоремы 2 вытекает 


$$ 


|u(x)|\le M C'_\rho(x)\sigma^2\exp(-\sigma x_3) 


|u_\sigma(x)|. 


$$ 


Согласно условию теоремы 3 и неравенствам (15), (16) имеем 


\begin{align*} 


|u_\sigma(x)|&\le \delta\int_S[|\Pi(y,x,\sigma)|+ 


|T(\partial_y,n)\Pi(y,x,\sigma)|]ds_y \le \\ 


% 


&\le \delta\sigma^2\exp(-\sigma x_3)\int_S[C_4(\rho)r^{-1} 


+C_5(\rho)r^{-2}]\exp(-\sigma y_3)ds_y \le \\ 


% 


&\le \delta\sigma^2\exp\sigma(h-x_3)\int_S[C_4(\rho)r^{-1}+ 


C_5(\rho)r^{-2}]ds_y. 


\end{align*} 


Теперь из (13) и (14) \ $|u(x)|\le C_\rho(x)\sigma^2\exp(-\sigma 


x_3)[M+\delta\exp\sigma h]$. 


Здесь $M=\delta\exp(\sigma h)$ или 


$\sigma=\frac{1}{h}\ln\big(\frac{M}{\delta}\big)$. В результате получим 


$$ 


|u(x)|\le C_\rho(x)\frac{2M^{1-x_3/h}}{h^2}\sigma^{x_3/h}\ln\bigg( 


\frac{M}{\delta}\bigg)^2.\quad\square 


$$ 


 


 


Теперь предположим, что вместо $u(y)$ и $T(\partial_y,n)u(y)$ на $S$ заданы 


их 


непрерывные приближения $f_\delta(y)$ и $g_\delta(y)$ соответственно, 


$$ 


\max_S|u(y)-f_\delta(y)|+\max_S|T(\partial_y,n)u(y)- 


g_\delta(y)|<1, 


$$ 


предполагаем, что $S$ удовлетворяет условиям Ляпунова. 


 


Обозначим 


\begin{equation} 


u_{\sigma\delta}(x)=\int_S[\Pi(y,x,\sigma)g_\delta(y)-f_\delta(y) 


\{T(\partial_y,n)\Pi(y,x,\sigma)\}]ds_y,\quad x\in D. 


\end{equation} 


 


 


\begin{thm} 


Пусть $u(y)\in A_\rho(D)$ удовлетворяет граничному условию 


$(7)$. Тогда 


$$ 


|u(x)-u_{\sigma\delta}(x)|\le C_\rho(x)\frac{2M^{1-x_3/h}}{h^2} 


\delta^{x_3/h}\ln\bigg(\frac{M}{\delta}\bigg)^2, 


$$ 


где $\sigma=h^{-1}\ln\big(\frac{M}{\delta}\big)$, 


$C_\rho(x)=C(\rho)\int\limits_{\partial D}[r^{-1}+r^{-2}]ds_y$. 


\end{thm} 


 


 


{\bf Доказательство.} Согласно (8) и (17) имеем 


\begin{gather*} 


u(x)-u_{\sigma\delta}(x)=\int_\Sigma[\Pi(y,x,\sigma) 


\{T(\partial_y,n)u(y)\}\{T(\partial_y,n)\Pi(y,x,\sigma)\}]ds_y+ \\ 


% 


+\int_\Sigma[\Pi(y,x,\sigma)\{T(\partial_y,n)u(y)-g_\delta(y)\} 


-(u(y)-f_\delta(y))\{T(\partial_y,n)\Pi(y,x,\sigma)\}]ds_y. 


\end{gather*} 


В силу теорем 2 и 3 получим 


\begin{gather*} 


|u(x)-u_{\sigma\delta}(x)|\le M\int_\Sigma[|\Pi(y,x,\sigma)|+ 


|T(\partial_y,n)\Pi(y,x,\sigma)|]ds_y+ \\ 


% 


+\delta\int_S[|\Pi(y,x,\sigma)|+|T(\partial_y,n)\Pi(y,x,\sigma)|]ds_y\le 


C_\rho(x)\frac{2M^{1-x_3/h}}{h^2}\delta^{x_3/h}\ln\bigg( 


\frac{M}{\delta}\bigg)^2.\quad\square 


\end{gather*} 


 


 


Из доказанных теорем получаем 


 


 


\begin{cor} 


Предельные равенства 


$$ 


\lim_{\sigma\to x}u_\sigma(x)=u(x),\quad 


\lim_{\delta\to0}u_{\sigma\delta}(x)=u(x) 


$$ 


выполняются равномерно на каждом компакте из $D$. 


\end{cor} 


 


 


Формула (9) дает в явном виде приближенное решение задачи (1)--(2) в 


точке $x\in D$, формула (17) представляет приближенное решение, когда 


данные Коши на $S$ заданы приближенно. Эти формулы основаны на 


постановке и методе анализа, предложенных в [2], [4]. 


\medskip 


 


{\bf 2.} Приведем аналогичные результаты для системы уравнений 


термоупругости. Пусть $x=(x_1,x_2,x_3)$, $y=(y_1,y_2,y_3)$ --- точки 


трехмерного 


евклидова пространства $E^3$, односвязная неограниченная область $D\subset 


E^3$ лежит внутри слоя $0<y_3<h$, $h=\pi/\rho$, $\rho>0$, граница ее состоит 


из плоскости $\Sigma:y_3=0$ и гладкой поверхности $S$, заданной уравнением 


$y_3=F(y_1,y_2)$ и удовлетворяющей условиям $0<F(y_1,y_2)\le h$, $|\Grad 


F(y_1,y_2)|\le\const<\infty$, $(y_1,y_2)\in E^2$, 


т.\,е. $\partial D=\Sigma\cup S$. Рассмотрим матричный 


дифференциальный оператор 


\begin{gather*} 


B(\partial_x,\omega)=\|B_{mj}(\partial_x,\omega)\|_{4\times 4}, \\ 


% 


B_{mj}(\partial_x,\omega)=(1-\delta_{m4})(1-\delta_{4j}) 


[\delta_{mj}\mu(\Delta(\partial_x)+q\omega^2/\mu)+ 


(\lambda+\mu)\partial^2/\partial x_m\partial x_j]- \\ 


% 


-\delta_{4j}(1-\delta_{m4})\gamma\partial/\partial x_m+ 


i\omega\eta\delta_{m4}(1-\delta_{4j})\partial/\partial x_j 


+\delta_{m4}\delta_{4j}(\Delta(\partial_x)+i\omega/\chi),\quad 


m,j=1,2,3,4, 


\end{gather*} 


где $\Delta(\partial_x)$ --- оператор Лапласа, 


$\lambda$, $\mu$, $q$ --- постоянные Ламе и плотность 


упругой среды соответственно, $\omega$ --- частота колебания, $\delta_{mj}$ 


--- символ 


Кронекера, $\gamma=(3\lambda+2\mu)\alpha$, $\alpha$ --- коэффициент линейного 


температурного расширения 


среды, $\chi$ --- коэффициент температуропроводности, $i$ --- мнимая единица, 


$k$ --- коэффициент теплопроводности среды, $\eta=\gamma\theta_0/k$, 


$\theta_0$ --- температура среды в 


недеформированном состоянии. Уравнение термоупругоколебательного 


состояния среды $D$ в компонентах смешения принимает вид ([15], с.\,12) 


\begin{equation} 


B(\partial_x,\omega)U=0,\tag{1$'$} 


\end{equation} 


где $U=(u_1,u_2,u_3,u_4)=(u,u_4)$ --- четырехмерный вектор, $u_1$, $u_2$, 


$u_3$, как и выше, 


обозначают компоненты смещения, $u_4=\theta$ --- отклонение температуры среды 


от температуры $\theta_0$. Решение $U$ системы ($1'$) в области $D$ назовем 


регулярным, если $U\in C^2(D)\cap C^1(\ov D)$. 


 


Обозначим через $A(D)$ пространство всех регулярных в $D$ решений 


системы ($1'$), а через $A_\rho(D)$ --- класс функции из $A(D)$, 


удовлетворяющих следующему условию роста: 


$$ 


A_\rho(D)=\{U(y)\in A(D):|U(y)|+|\Grad U(y)|\le\exp 


(0(\exp\rho|y_1|)),\ \; y\to\infty,\ \; \rho>0\}. 


$$ 


 


Далее будем использовать матричный дифференциальный 


оператор термонапряжения 


$$ 


R(\partial_y,n(y))= 


\begin{Vmatrix} 


\ & T & \ & -\gamma n_1\\ 


\ & \ & \ & -\gamma n_2\\ 


\ & \ & \ & -\gamma n_2 \\ 


0 & 0 & 0 & -\partial/\partial n 


\end{Vmatrix} 


, 


$$ 


где $T$ означает оператор напряжения 


$$ 


T(\partial_y,n(y))=\|T_{mj}\|_{3\times 3}=\|\lambda n_m 


\partial/\partial y_j+\mu n_j \partial/\partial y_m+ 


\mu\delta_{mj}\partial/\partial n\|_{3\times 3}, 


$$ 


$n(y)=(n_1,n_2,n_3)$ --- единичный вектор внешней нормали к поверхности 


$\partial D$ в 


точке $y$. Через $\wt R(\partial_y,n(y))$ обозначим матрицу, которая 


получается из 


$R(\partial_y,n(y))$ заменой в последнем столбце $\gamma$ на $i\omega\eta$. 


\medskip 


 


{\it Постановка задачи}. Требуется определить регулярное решение $U$ 


системы ($1'$) в области $D$, исходя из ее данных Коши, заданных на 


поверхности $S$, 


\begin{equation} 


U(y)|_S=f(y),\quad R(\partial_y,n(y))U(y)|_S=g(y), \tag{2$'$} 


\end{equation} 


$f=(f_1,f_2,f_3,f_4)$, $g=(g_1,g_2,g_3,g_4)$ --- заданные непрерывные 


вектор-функции. 


 


Данная задача также является некорректной. Характер некорректности 


такой же, как у задачи Коши для уравнения Гельмгольца. Для случая 


$\omega=0$, $\gamma=0$ рассматриваемая задача для специальных неограниченных 


классов областей исследована в [10]--[14]. 


 


Следуя [2], приведем 


 


\begin{defn} 


Матрицей Карлемана задачи ($1'$)--($2'$) называется $4\times4$-матрица 


$\Pi(y,x,\omega,\gamma,\sigma)$, удовлетворяющая следующим двум условиям: 


 


1) $\Pi(y,x,\omega,\gamma,\sigma)=\Gamma(y,x,\omega,\gamma)+ 


G(y,x,\omega,\gamma,\sigma)$, 


 


\noindent{}где $\sigma$ --- положительный числовой параметр, матрица 


$G(y,x,\omega,\gamma,\sigma)$ по 


переменной $y$ удовлетворяет системе ($1'$) всюду в области $D$, 


$\Gamma(y,x,\omega,\gamma)$ --- 


матрица фундаментальных решений уравнений ($1'$) (см. [15], с.\,38); 


 


2) $\int\limits_\Sigma|\Pi|+|(\wt R\Pi^*)^*|ds_y\le\varepsilon(\sigma)$, 


 


\noindent{}где $\varepsilon(\sigma)\to0$ при $(\sigma)\to\infty$; здесь и 


далее $\Pi^*$ означает матрицу, 


сопряженную матрице $\Pi$, а $|\Pi|$ --- евклидова норма матрицы $\Pi$ , 


т.\,е. 


$|\Pi|=\Big(\sum\limits^4_{m,j=1}\Pi^2_{mj}\Big)^{1/2}$. 


\end{defn}





 


С целью построения приближенного решения задачи ($1'$) --($2'$) 


рассмотрим матрицу 


\begin{gather} 


\Pi(y,x,\omega,\gamma,\sigma)=\|\Pi_{mj}(y,x,\omega,\gamma,\sigma) 


\|_{4\times 4}, \\ 


% 


\Pi_{mj}(y,x,\omega,\gamma,\sigma)= 


\frac{(1-\delta_{m4})(1-\delta_{j4})}{2\pi}\bigg[\frac{\delta_{mj}}{\mu} 


\Phi_\sigma(y,x,i\lambda_k)-\sum^3_{k=1}\alpha_k 


\frac{\partial^2}{\partial x_m\partial x_j}\Phi_\sigma 


(y,x,i\lambda_k)\bigg] + \notag \\ 


% 


+\frac{i\omega\eta}{2\pi}\delta_{m4}(1-\delta_{j4})\sum^3_{k=1} 


\beta_k\frac{\partial}{\partial x_j}\Phi_\sigma(y,x,i\lambda_k)-\notag \\ 


% 


-\frac{\gamma}{2\pi}\delta_{i4}(1-\delta_{m4})\sum^3_{k=1}\beta_k 


\frac{\partial}{\partial x}\Phi_\sigma(y,x,i\lambda_k)+ 


\sum^3_{k=1}\frac{\delta_{j4}\delta_{m4}}{2\pi}\gamma_k 


\Phi_\sigma(y,x,i\lambda_k),\notag 


\end{gather} 


где постоянные $\alpha_k$, $\beta_k$, $\gamma_k$, $\lambda_k$ явно выражаются 


через постоянные 


термоупругости --- коэффициенты системы ($1'$) ([15], с.\,432) 


\begin{gather} 


\Phi_\sigma(y,x,\Lambda)=\frac{1}{2\pi}\int^\infty_0\im 


\frac{K(\sigma(\omega-x_3))}{\omega-x_3} 


\frac{\cos\Lambda u\,du}{\sqrt{u^2+r^2_1}},\notag \\ 


% 


K(\sigma\omega)=(\omega-x_3+2h)^{-1}\exp(\sigma\omega+\omega^2),\quad 


\omega=i\sqrt{u^2+r^2_1}+y_3, \\ 


% 


r^2_1=(y_1-x_1)^2+(y_2-x_2)^2,\quad r_1>0,\quad 0<x_3<h.\notag 


\end{gather} 


 


Из результатов работы [17] вытекает 


 


 


\begin{lem} 


Функция $\Phi_\sigma(y,x,\Lambda)$ является функцией Карлемана для 


уравнения Гельмгольца, т.\,е. обладает следующими двумя свойствами\/${:}$ 


$$ 


\Phi_\sigma(y,x,\Lambda)=\frac{\exp(\Lambda r)}{4\pi r} 


+\nu(y,x,\Lambda,\sigma),\quad r=|x-y|, 


$$ 


где $\nu(y,x,\Lambda,\sigma)$ --- некоторая функция, определенная для всех 


значений $y$, $x$ и 


удовлетворяющая условию Гельмгольца 


\begin{gather*} 


\Delta(\partial_y)\nu-\Lambda^2\nu=0,\quad y\in D,\\ 


% 


\int_\Sigma\bigg(|\Phi_\sigma|+\bigg| 


\frac{\partial\Phi(\sigma)}{\partial n}\bigg|\bigg)ds_y\le 


C(\Lambda,D)\sigma\exp(-\sigma x_3), 


\end{gather*} 


где $C(\Lambda,D)$ --- постоянная. 


\end{lem} 


 


 


Верна аналогичная лемма для системы ($1'$). 


 


 


\begin{lem} 


Матрица $\Pi(y,x,\omega,\gamma,\sigma)$, определенная равенствами $(18)$ и 


$(19)$, является матрицей Карлемана для задачи $(1')$--$(2')$. 


\end{lem} 


 


 


Доказательства лемм 3, 4 аналогичны доказательствам лемм 1, 2. 


Положим 


\begin{equation} 


2U^\sigma(x)=\int_S(\Pi(y,x,\omega,\gamma,\sigma)g(y)- 


(\wt R(\partial_y,n(y))\Pi^*(y,x,\omega,\gamma,\sigma)) 


f(y))ds_y,\quad x\in D. 


\end{equation} 


 


Имеет место 


 


 


\begin{thm} 


Пусть $U\in A_\rho(D)$ удовлетворяет условию 


\begin{align} 


|U(y)|+|R(\partial_y,n)U(y)|&\le1,\quad y\in\Sigma,\notag \\ 


% 


|U(x)|+|R(\partial_y,n)U(x)|&\le\exp(o(\exp\rho\sqrt{x^2_1+x^2_2})), 


\quad x\to\infty,\ \ x\in D. 


\end{align} 


Тогда при $\sigma\ge1$ справедливо неравенство 


$$ 


|U(x)-U^\sigma(x)|\le C_1(\lambda,\nu,\omega,\gamma,D) 


\sigma^3\exp(-\sigma x_3),\quad x\in D. 


$$ 


\end{thm} 


 


 


Заметим, что в теореме 5 при сделанных предположениях верна 


интегральная формула 


$$ 


U(x)=\int_{\partial D}(\Pi R(\partial/\partial y,n)U(y)- 


\wt R(\partial/\partial y,n)\Pi^*)^* U(y)ds,\quad x\in D. 


$$ 


 


Приведем результат, который позволяет вычислить $U(x)$ приближенно, 


когда на поверхности $S$ вместо $U(y)$ и 


$R(\partial/\partial y,n(y))U(y)$ заданы их 


непрерывные приближения $f_\delta(y)$ и $g_\delta(y)$, для которых 


$$ 


\max_S|f(y)-f_\delta(y)|+\max_S|g(y)-g_\delta(y)| 


\le\delta,\quad 0<\delta<1. 


$$ 


Предполагаем, что $S$ удовлетворяет условиям Ляпунова. Определим 


функцию 


\begin{equation} 


2U^{\sigma\delta}(x)=\int_S\bigg(\Pi g_\delta(y)-\wt R\bigg( 


\frac{\partial}{\partial y},n\bigg)\Pi^*\bigg)f_\delta(y)ds_y,\quad x\in D, 


\end{equation} 


где $\sigma=\frac{1}{x^0_3}\ln\frac{1}{\delta}$, $x^0_3=\max\limits_D x_3$. 


Тогда имеет место 


 


 


\begin{thm} 


Пусть $U\in A_\rho(D)$ удовлетворяет граничному условию 


$(21)$ на всей границе. Тогда справедливо неравенство 


$$ 


|U(x)-U^{\sigma\delta}(x)|\le C_2(\lambda,\mu,\omega,\gamma,D) 


\delta^{x_3/x^0_3}\bigg(\ln\frac{1}{\delta}\bigg)^3,\quad x\in D. 


$$ 


\end{thm} 


 


 


Доказательства теорем 5, 6 аналогичны доказательствам теорем 2, 4. 


 


 


\begin{cor} 


Предельные равенства 


$$ 


\lim_{\sigma\to\infty}U^\sigma(x)=U(x),\quad 


\lim_{\delta\to0}U^{\sigma\delta}(x)=U(x), 


$$ 


выполняются равномерно на каждом компакте из $D$. 


\end{cor} 


 


 


Формула (19) дает в явном виде приближенное решение задачи $(1')$--$(2')$ 


в точке $x\in D$, а формула (21) изображает приближенное решение, когда 


данные Коши на $S$ заданы приближенно. Эти формулы получены на 


основе результатов М.М.\,Лаврентьева (см. [2], [4]). 


 


В заключение автор выражает признательность профессору 


Ш.Я.\,Ярмухамедову за постановку задачи и обсуждения в процессе ее 


решения. 


 


 


\begin{thebibliography}{99} 


 


\bibitem{1} 


Адамар Ж. {\em Задача Коши для линейных уравнений с частными 


производными гиперболического типа}. -- М.: Наука, 1978. -- С.\,38--70. 


\bibitem{2} 


Лаврентьев М.М. {\em О некоторых некорректных задачах 


математической физики}. -- Новосибирск: ВЦ СО АН СССР, 1962. -- 92~с. 


\bibitem{3} 


Берс А., Джон Ф., Шехтер М. {\em Уравнения с частными производными}. -- 


М.: Мир, 1966. -- 351~с. 


\bibitem{4} 


Лаврентьев М.М. {\em О задаче Коши для линейных эллиптических уравнений 


второго порядка} // ДАН СССР. -- 1957. -- T.\,112. -- \No\,2. -- 


С.\,195--197. 


\bibitem{5} 


Мергелян С.Н. {\em Гармоническая аппроксимация и приближенное решение 


задачи Коши для уравнения Лапласа} // УМН. -- 1956. -- T.\,11. -- 


Вып.\,5. -- С.\,3--26. 


\bibitem{6} 


Иванов В.К. {\em Задача Коши для уравнения Лапласа в бесконечной 


полосе} // Дифференц. уравнения. -- 1965. -- Т.1. -- \No\,1. -- С.\,131--136. 


\bibitem{7} 


Ярмухамедов Ш.Я. {\em О задаче Коши для уравнения Лапласа} // ДАН 


СССР. -- 1977. -- Т.\,235. -- \No\,2. -- С.\,281--283. 


\bibitem{8} 


Ярмухамедов Ш.Я., Ишанкулов Т.И., Махмудов О.И. {\em О задаче Коши 


для системы уравнений теории упругости в пространстве} // Сиб. матем. журн. 


-- 1992. -- Т.\,33. -- \No\,2. -- С.\,281--283. 


\bibitem{9} 


Айзенберг Л.А., Тарханов Н.Н. {\em Абстрактная формула Карлемана} // 


ДАН СССР. -- 1988. -- Т.\,298. -- \No\,6. -- С.\,1292--1296. 


\bibitem{10} 


Махмудов О.И. {\em Задача Коши для системы уравнений 


пространственной теории упругости в перемещениях} // Изв. вузов. 


Математика. -- 1994. -- \No\,1. -- С.\,54--61. 


\bibitem{11} 


Махмудов О.И., Ни\"езов И.Э. {\em Регуляризация решения задачи Коши 


для системы уравнений теории упругости в перемещениях} // Сиб. матем. 


журн. -- 1998. -- Т.\,39. -- \No\,1. -- С.\,27--35. 


\bibitem{12} 


Ни\"езов И.Э. {\em Задача Коши для системы теории упругости на 


плоскости} // Узб. матем. журн. -- 1996. -- \No\,1. -- С.\,44--51. 


\bibitem{13} 


Ишанкулов Т.И., Махмудов О.И. {\em Задача Коши для системы 


уравнений термоупругости в пространстве} // Изв. вузов. Математика. -- 


1999. -- \No\,6. -- С.\,27--32. 


\bibitem{14} 


Махмудов О.И., Ни\"езов И.Э. {\em Регуляризация решения задачи Коши 


для системы теории упругости} // Матем. заметки. -- 2000. -- Т.\,68. -- 


\No\,4. -- С.\,548--554. 


\bibitem{15} 


Купрадзе В.Д., Гегелиа Т.Г., Башевейшвили М.О., Бурчуладзе Т.В. 


{\em Трехмерные задачи математической теории упругости и 


термоупругости. Классическая и микрополярная теория. Статика, гармонические 


колебания, динамика. Основы и методы решения}. -- 2-е изд. -- 


М.: Наука, 1976. -- 663~с. 


\bibitem{16} 


Тихонов А.Н. {\em О решении некорректно поставленных задач и методе 


регуляризации} // ДАН СССР. -- 1966. -- T.\,151. -- \No\,3. 


-- С.\,501--504. 


\bibitem{17} 


Тарханов Н.Н. {\em О матрице Карлемана для эллиптических систем} // 


ДАН СССР. -- 1985. -- Т.\,284. -- \No\,2. -- С.\,294--297. 


 


\end{thebibliography} 


 


\vskip 2cm 


 


\postup{\it Самаркандский государственный}{\qquad \it Поступили} 


\postup{\it университет}{\it первый вариант $04.06.2001$} 


\postup{}{\it окончательный вариант $22.03.2002$} 


 


\label{end} 


\end{document} 


